% =====================================================================
%  §3 Calabi–Yau manifolds in the large complex structure limit
%  原文 p13–p21
% =====================================================================
\section{大复结构极限下的 Calabi--Yau 流形}

\subsection{Calabi--Yau 流形的极大退化}

设 $\Cmer=\{f=\sum_{n\ge n_0}a_nq^n\}$ 为单复变量亚纯函数在 $q=0$ 处的芽之域。

设 $\mathcal{X}_{mer}$ 是 $\Cmer$ 上的代数 $n$ 维 Calabi--Yau 流形（即：$\mathcal{X}_{mer}$
是 $\Cmer$ 上的光滑射影流形，典范类平凡，$K_{\mathcal{X}}=0$）。我们固定一个代数处处非零的
体积元 $vol\in\Gamma(\mathcal{X}_{mer},K_{\mathcal{X}})$。配对 $(\mathcal{X}_{mer},vol)$
定义了一族单参数复 Calabi--Yau 流形 $(X_q,vol_q)$，$0<|q|<r_0$，其中 $r_0>0$。

设 $[\omega]\in H^2_{\DR}(\mathcal{X}_{mer})$ 是丰富锥中的上同调类。于是对每个满足
$0<|q|<r_0$ 的 $q$，它都在 $X_q$ 上定义了一个凯勒类 $\omega_q$。由 Yau 定理，存在 $X_q$ 上
唯一的 Calabi--Yau 度量 $g_{X_q}$，其凯勒类为 $[\omega_q]$。

由奇点消解可以推出，当 $q\to0$ 时有如下公式：
\[
  \int_{X_q}vol_q\wedge\overline{vol_q}
  =C(\log|q|)^m|q|^{2k}\bigl(1+o(1)\bigr),
\]
其中 $C\in\C^*$，$k\in\Z$，$0\le m\le n=\dim(\mathcal{X}_{mer})$。

\begin{dfn}
若上述公式中 $m=n$，则称 $\mathcal{X}_{mer}$ 在 $q=0$ 处有极大退化（maximal degeneration）。
\end{dfn}

\begin{yizhu}[极大退化的含义]
体积 $\int_{X_q}vol_q\wedge\overline{vol_q}$ 随 $|q|\to0$ 按 $|q|^{2k}(\log|q|)^m$ 增长：
$|q|^{2k}$ 是幂次因子，而对数因子 $(\log|q|)^m$ 来自周期积分的单值性，其幂次至多为
$n=\dim X$。取到上界 $m=n$ 就是「极大」一词的确切含义；紧随其后的命题 1 把它等价地刻画为
单值算子 $T$ 的幂零指数恰为 $n+1$。
\end{yizhu}

下面我们说明，这个定义与通常用 Hodge 结构的形变给出的定义等价（见 \cite{Mo}）。

\begin{prp}
Calabi--Yau 流形 $\mathcal{X}_{mer}$ 有极大退化，当且仅当：对所有充分小的 $q$，存在向量
$v\in H^n(X_q,\C)$，使得
\[
  (T-\Id)^{n+1}v=0,\qquad (T-\Id)^{n}v\ne0,
\]
其中 $T$ 是单值算子。
\end{prp}

\begin{proof}
首先注意，体积 $\int_{X_q}vol_q\wedge\overline{vol_q}$ 可以在上同调中计算，即等于中间上同调
$H^n(X_q,\C)$ 的（本原部分中的）Poincaré 配对
$\langle[vol_q],[\overline{vol_q}]\rangle$。

我们可以假设（在添入 $q$ 的一个方根、过渡到覆盖之后）算子 $T$ 是幂幺的。我们用乘以
\[
  q^{-\log(T)/2\pi i}
  =\sum_{k=0}^{n}\left(\frac{\log(q)}{-2\pi i}\right)^{k}\frac{(\log T)^{k}}{k!}
\]
的方式，把穿孔圆盘上的丛 $H^n(X_q,\C)$ 平凡化。

幂零轨道定理（nilpotent orbit theorem）说，Hodge 滤过延拓为平凡化丛在整个圆盘（包括
$q=0$ 这一点）上的全纯滤过。因此，丛 $q^{-\log(T)/2\pi i}(H^{n,0})$ 延拓为圆盘上的线丛，
而截面 $q^{-\log(T)/2\pi i}([vol_q])$ 是这个线丛的非零亚纯截面。乘以 $q$ 的适当幂次之后，
我们可以进一步假设该截面全纯且在 $q=0$ 处非零。我们把这个全纯截面记作 $a(q)$，$a(0)\ne0$。

现在来计算体积：
\begin{align*}
  \langle[vol_q],[\overline{vol_q}]\rangle
  &=\sum_{k,l\ge0}\left\langle
     \left(\frac{\log(q)}{2\pi i}\right)^{k}\frac{(\log T)^{k}}{k!}a(q),\;
     \left(\frac{\log(\bar q)}{-2\pi i}\right)^{l}\frac{(\log T)^{l}}{l!}\overline{a(q)}
   \right\rangle\\[2pt]
  &=\sum_{k,l\ge0}\left(\frac{\log(q)}{2\pi i}\right)^{k}
     \left(\frac{\log(\bar q)}{-2\pi i}\right)^{l}
     \frac{(-1)^{l}}{k!\,l!}
     \bigl\langle(\log T)^{k+l}a(q),\overline{a(q)}\bigr\rangle .
\end{align*}
这里我们用到：算子 $\log(T)$ 是实的，且关于 Poincaré 配对是斜对称的。由等式
$\log(T)^{n+1}=0$（它自动成立）可知，上式中所有 $k+l>n$ 的项都消失。$k+l=n$ 的项之贡献等于
\[
  \frac{(\log|q|)^{n}}{(\pi i)^{n}n!}
  \bigl\langle(\log T)^{n}a(q),\overline{a(q)}\bigr\rangle .
\]
由极化 Hodge 结构的形变的标准性质（例如见 \cite{De}）容易推出：若 $\log(T)^{n}\ne0$，
则事实上 $\log(T)^{n}(a(0))\ne0$，并且
$\bigl\langle(\log T)^{n}a(0),\overline{a(0)}\bigr\rangle\ne0$。于是命题的陈述变得显然。
\end{proof}

注意，\cite{De} 中施加了一个稍强的极大退化条件：$H^*(X_q,\C)$ 上由单值算子给出的权滤过
应当与 Hodge 滤过互补。

我们回忆 Gromov--Hausdorff 度量 $\rho_{GH}$ 的定义。它是紧度量空间的等距类所成空间上的
一个度量。若存在度量空间 $M$，它以 $M_1$、$M_2$ 为度量子空间同时包含二者，使得 $M_1$
落在 $M_2$ 的 $\eps$-邻域之内、且反之亦然，则称两个度量空间 $M_1$ 与 $M_2$ 在 $\rho_{GH}$
下 $\eps$-接近。此时 $\rho_{GH}(M_1,M_2)$ 取为这种 $\eps$ 的最小值。

我们把 Calabi--Yau 度量重新标度：$g^{new}_{X_q}=g_{X_q}/\operatorname{diam}(X_q,g_{X_q})^2$。
这样就得到一族直径为 $1$ 的黎曼流形 $X_q^{new}=(X_q,g^{new}_{X_q})$，以 $q$ 为参数。

\begin{conj}
若 $\mathcal{X}_{mer}$ 在 $q=0$ 处有极大退化，则 $X_q^{new}$ 在 Gromov--Hausdorff
度量下存在极限 $(\bar Y,g_{\bar Y})$，使得：
\begin{enumerate}[label=\textup{\alph*}),leftmargin=2em,itemsep=2pt]
  \item $(\bar Y,g_{\bar Y})$ 是紧度量空间，它包含一个 $n$ 维光滑定向黎曼流形 $(Y,g_Y)$
        作为稠开的度量子空间。$Y^{sing}=\bar Y\setminus Y$ 的 Hausdorff 维数不超过 $n-2$。
  \item $Y$ 带有整仿射结构（integral affine structure）。这是指：$Y$ 上带有一个无挠平坦
        联络 $\nabla$，其和乐含于 $SL(n,\Z)$ 之中。
  \item 度量 $g_Y$ 有势。这是指：在仿射坐标下它局部由对称矩阵
        $(g_{ij})=(\partial^2K/\partial x_i\partial x_j)$ 给出，其中 $K$ 是光滑函数
        （定义到相差一个仿射函数，即一个线性函数与一个常数之和）。
  \item 在仿射坐标下度量的体积元为常数，
        $\det(g_{ij})=\det(\partial^2K/\partial x_i\partial x_j)=\mathrm{const}$
        （实 Monge--Ampère 方程）。
\end{enumerate}
\end{conj}

本节末尾我们将基于微分几何的考虑，给出猜想的一个不严格的解释。

\begin{rmk}
1) 由于度量 $g_Y$ 决定的矩阵 $(g_{ij})$ 是正定的，函数 $K$ 是凸的。特别地，$K$ 的
Legendre 变换局部有定义。这一事实将在后面讨论 Monge--Ampère 流形的对偶时用到。

2) 看起来可信的是：当所有 $X_q$ 都单连通、且 $h^{k,0}(X_q)=0$（$0<k<n$）时，度量空间
$\bar Y$ 是 $n$ 维同调球。在所有例子中它事实上同胚于 $S^n$。
\end{rmk}

该猜想为下述紧化开辟了道路：给定 Calabi--Yau 流形 $M$ 上 Calabi--Yau 度量的模空间，
可以把一切满足「对某个 $q'$ 有 $X_{q'}=M$」的单参数极大退化 $X_q$ 所对应的配对
$(\bar Y,\R^*_+\cdot g_{\bar Y})$ 之集合，作为边界分支添入。这对应于 Calabi--Yau
流形模空间中一个「尖点」的选取。这种选取通常用某些代数几何数据来描述：单值算子的作用、
Hodge 结构的形变、特殊纤维的混合 Hodge 结构，等等。前一个猜想给出了尖点的一种纯粹
「度量」的描述。

由猜想的 b) 部分可知，可以选取 $\nabla$-协变格 $T^{\Z}_{Y,y}\subset T_{Y,y}$，$y\in Y$。
设给定满足猜想 a)--c) 的三元组 $(Y,g_Y,\nabla)$，并在切丛 $T_Y$ 中固定了一个协变格
$T^{\Z}_Y$。那么我们可以构造一族\emph{非紧}复 Calabi--Yau 流形，带 Ricci 平坦的凯勒度量。
具体地说，设 $X^{\eps}$ 是环面丛 $p_\eps:X^\eps\to Y$ 的全空间，其纤维为
$T_{Y,y}/\eps T^{\Z}_{Y,y}$，$y\in Y$，$0<\eps\le\eps_0$。切丛 $T_Y$ 的全空间 $TY$
带有典范的复结构，它来自同构
\[
  T_{TY}\simeq\pi^*T_Y\oplus\pi^*T_Y\simeq\pi^*T_Y\otimes\C,
\]
其中 $\pi:TY\to Y$ 是典范投影（这里我们用到 $Y$ 上的仿射结构）。利用同一等同，我们在
$TY$ 上引入度量 $g_{TY}=\pi^*g_Y\oplus\pi^*g_Y$。容易看出，$g_{TY}$ 是以 $\pi^*K$ 为势的
凯勒度量。由 Monge--Ampère 方程可知，度量 $g_{TY}$ 是 Ricci 平坦的。过渡到商，我们在
$X^\eps$ 上得到复结构 $J_{X^\eps}$ 与 Ricci 平坦的凯勒度量 $g_{X^\eps}$。

设 $U\subset Y$ 是单连通开子集。则环面 $T^n\simeq T_{Y,y}/T^{\Z}_{Y,y}$ 作用在
$p_\eps^{-1}(U)$ 上，$y\in U$（对不同的点 $y\in U$，借助联络 $\nabla$ 把这些环面互相
等同）。由此可知，对任意 $t\in H^1\bigl(Y,(T_Y/T^{\Z}_Y)_{discr}\bigr)$（系数取在把环面
看作抽象群所得的局部系中），都可以定义一个扭流形 $X^{\eps,t}$，它是环面纤维化
$p_{\eps,t}:X^{\eps,t}\to Y$ 的全空间。

粗略地说，下一个猜想断言：Calabi--Yau 流形族 $X_q^{new}$（$q=e^{-1/\eps}$）在极大退化点
$\eps=0$ 附近的「首项渐近」，在差一个扭的意义下同构于上面描述的环面丛所对应的族
$(X^\eps,J_{X^\eps})$。

更精确地说，我们把它表述如下。

\begin{conj}
设 $(\mathcal{X}_{mer},vol)=(X_q,vol_q)$ 是一族单参数极大退化 Calabi--Yau 流形，
$X_q^{new}$ 是如前重新标度度量之后所得的族。则存在常数 $C>0$ 与函数 $t(q)$，使得凯勒流形
$X_q^{new}$ 与 $X^{\eps(q),t(q)}$（其中 $\eps(q)=C(\log|q|)^{-1}$）在 $q\to0$ 时按如下
意义互相接近：

对任意 $\delta>0$，存在分解 $X_q=X_q^{sm}\sqcup X_q^{sing}$ 与光滑流形的嵌入
\[
  j_q:X_q\to p_{\eps(q),t(q)}^{-1}\bigl(Y\setminus(Y^{sing})^{\delta}\bigr),
\]
其中 $(Y^{sing})^{\delta}$ 是 $Y^{sing}$ 的 $\delta$-邻域，使得：
\begin{enumerate}[label=\textup{\alph*}),leftmargin=2em,itemsep=2pt]
  \item $(X_q,X_q^{sing})$ 在 Gromov--Hausdorff 度量下收敛到配对
        $(\bar Y,Y^{sing})$。
  \item $j_q$ 把切空间 $T_xX_q$ 与 $T_{j_q(x)}X^{\eps(q),t(q)}$ 上的内积与复结构等同起来，
        误差为 $o(1)$，且对 $x\in X_q^{sm}$ 一致。
\end{enumerate}
\end{conj}

猜想 1 与猜想 2 有如下动机。一般而言，对于 Ricci 曲率非负的退化黎曼度量族，人们期待
用带奇点的纤维化塔（坍缩）来描述（与 §2.3 比较）\footnote{度量空间紧化纲领中的若干步骤
现在已经完成（例如见 \cite{CC}），但仍有许多未澄清的问题。}。在凯勒流形的情形，简单坍缩有
两种基本图景。第一种情形是底空间与纤维都是凯勒流形。第二种情形是纤维为 $m$ 维平坦全实
环面，而底空间局部看起来是 $\R^m$ 中一个区域与某个凯勒流形的乘积。体积渐近行为中的对数
因子只应来自环面纤维。因此，对数可能出现的最高幂次，只有当存在纯粹环面纤维化之塔时才会
出现。看起来，凯勒类的固定（精确到一个标量）禁止了多重坍缩。这些考虑给出了我们的猜想的
一个直觉「解释」。

\begin{rmk}
在本文写作期间我们获悉，M.~Gross 与 P.~Wilson 独立提出了与我们的猜想类似的猜想
（见 \cite{GW}）。\cite{GW} 的一项显著成果，是在退化 K3 曲面的情形验证了这些猜想，并精确
描述了奇纤维附近度量的行为。此外，最近一篇预印本 \cite{Le} 从类似的观点讨论了镜像对称。
本文正文将考虑复阿贝尔簇的退化；在这种情形，猜想显然成立。
\end{rmk}

\subsection{Monge--Ampère 流形与环面纤维化的对偶}

本节我们为几何的镜像对称——即被理解为环面纤维化之间的对偶——提供一套数学语言。

\begin{dfn}
Monge--Ampère 流形（Monge--Ampère manifold）是一个三元组 $(Y,g,\nabla)$，其中 $(Y,g)$
是带度量 $g$ 的光滑黎曼流形，$\nabla$ 是 $T_Y$ 上的平坦联络，且满足：
\begin{enumerate}[label=\textup{\alph*}),leftmargin=2em,itemsep=2pt]
  \item $\nabla$ 在 $Y$ 上定义了一个仿射结构；
  \item 在仿射坐标 $(x_1,\dots,x_n)$ 下，$g$ 的矩阵 $\bigl((g_{ij})\bigr)$ 由
        $\bigl((g_{ij})\bigr)=\bigl((\partial^2K/\partial x_i\partial x_j)\bigr)$ 给出，
        其中 $K$ 是某个光滑实值函数；
  \item 满足 Monge--Ampère 方程
        $\det\bigl((\partial^2K/\partial x_i\partial x_j)\bigr)=\mathrm{const}$。
\end{enumerate}
\end{dfn}

Monge--Ampère 流形曾在 \cite{CY} 中以另一个名称被研究过，那里证明了：若 $Y$ \emph{紧致}，
则它的有限覆盖是环面。

我们来看一个由 K3 曲面的镜像对称所启发的（非紧）例子（另见 \cite{GW}）。设 $S$ 是带全纯
处处非零体积形式 $vol_S$ 的复曲面，$\pi:S\to C$ 是到复曲线 $C$ 上的全纯纤维化，其纤维为
非奇异椭圆曲线。

我们在 $C$ 上定义度量 $g_C$，取为与 $(1,1)$-形式 $\pi_*(vol_S\wedge\overline{vol_S})$
相伴的凯勒度量。在 $C$ 上（局部地）选取 $H_1(\pi^{-1}(x),\Z)$ 的一组基 $(\gamma_1,\gamma_2)$，
$x\in C$。用公式
\[
  \alpha_i=\mathrm{Re}\Bigl(\int_{\gamma_i}vol_S\Bigr),\qquad i=1,2
\]
在 $C$ 上定义两个闭 $1$-形式。

由此 $\alpha_i=dx_i$，其中 $x_i$（$i=1,2$）是某些函数。我们称 $(x_1,x_2)$ 为仿射坐标，
便在 $C$ 上定义了一个仿射结构及相应的联络 $\nabla$。可以直接验证 $(C,g_C,\nabla)$ 是
Monge--Ampère 流形。在 K3 曲面的椭圆纤维化的典型例子中，得到
$C=\C P^1\setminus\{z_1,\dots,z_{24}\}$，其中 $\{z_1,\dots,z_{24}\}$ 是 $\C P^1$ 中
$24$ 个互不相同的点。

回到一般情形。我们可以把我们猜想的一部分重新表述为：极大退化 Calabi--Yau 流形族的
Gromov--Hausdorff 极限的光滑部分，是带整仿射结构的 Monge--Ampère 流形。

Monge--Ampère 方程的局部解上有一个众所周知的对偶。

\begin{lmm}
设 $U\subset\R^n$ 是 $\R^n$ 中的凸开区域，配有标准仿射坐标 $(x_1,\dots,x_n)$，
$K:U\to\R$ 是满足 Monge--Ampère 方程的凸函数。则 Legendre 变换
\[
  \widehat K(y_1,\dots,y_n)
  =\max_{x\in U}\Bigl(\sum_i x_iy_i-K(x_1,\dots,x_n)\Bigr)
\]
也满足 Monge--Ampère 方程。
\end{lmm}

\begin{proof}
$L=dK$ 的图像是 $T^*\R^n=\R^n\oplus(\R^n)^*$ 中的拉格朗日子流形。设 $p_1$、$p_2$ 是到
两个直和项的典范投影。它们是局部微分同胚。由于 $K$ 定义到相差一个仿射函数，图像本身定义到
相差平移。Monge--Ampère 方程对应于条件
$p_1^*(vol_{\R^n})=p_2^*(vol_{\R^{*n}})$，其中 $vol_{\R^n}$（相应地 $vol_{\R^{*n}}$）
表示 $\R^n$（相应地 $\R^{*n}$）上的标准体积形式。流形 $L$ 可以看作 $d\widehat K$ 的图像。
因此 $\widehat K$ 也满足 Monge--Ampère 方程。引理证毕。
\end{proof}

流形 $L$ 带有由 $\R^n\oplus(\R^n)^*$ 上的不定度量 $\sum_i dx_idy_i$ 诱导的黎曼度量
$g_L$。在坐标 $(x_1,\dots,x_n)$ 下，这个度量由矩阵
$(\partial^2K/\partial x_i\partial x_j)$ 给出；在对偶坐标下，它由矩阵
$(\partial^2\widehat K/\partial y_i\partial y_j)$ 给出。于是 $L$ 上有一个度量与两个仿射
结构（坐标空间上标准仿射结构的拉回）。因此 $L$ 上有两个 Monge--Ampère 流形结构。容易
看出，这些局部图景可以拼贴起来。由此得到下述结果。

\begin{prp}
对给定的 Monge--Ampère 流形 $(Y,g_Y,\nabla_Y)$，存在典范定义的对偶 Monge--Ampère 流形
$(Y^{\vee},g^{\vee}_Y,\nabla^{\vee}_Y)$，使得 $(Y,g_Y)$ 与 $(Y^{\vee},g^{\vee}_Y)$ 作为
黎曼流形等同，且局部系 $(T_{Y^{\vee}},\nabla^{\vee}_Y)$ 自然同构于与 $(T_Y,\nabla_Y)$
对偶的局部系。
\end{prp}

\begin{crl}
若 $\nabla_Y$ 在 $Y$ 上定义整仿射结构（即 $\nabla_Y$ 的和乐含于 $GL(n,\Z)$ 中），则
$\nabla^{\vee}_Y$ 在 $Y^{\vee}$ 上定义整仿射结构。作为对偶协变格，取关于度量 $g_Y$ 与
$T^{\Z}_Y$ 对偶的格 $(T^{\Z}_Y)^{\vee}$。
\end{crl}

\begin{yizhu}[「对偶」改变的是什么]
命题 2 中的对偶并不改变底流形与度量：$(Y^{\vee},g^{\vee}_Y)$ 与 $(Y,g_Y)$ 作为黎曼流形
等同，改变的只是平坦联络——借度量把 $T_Y$ 与 $T^*_Y$ 等同之后，$\nabla^{\vee}_Y$ 就是
$\nabla_Y$ 的对偶联络。因此推论 1 中的对偶协变格 $(T^{\Z}_Y)^{\vee}$ 是 $T^{\Z}_Y$ 在
度量对偶下的像，两个整仿射结构互为对偶；这正是本节所谓「环面纤维化之对偶」的来源：纤维
$T_{Y,y}/\eps T^{\Z}_{Y,y}$ 的对偶纤维是 $T^*_{Y,y}/\eps(T^{\Z}_Y)^{\vee}_y$。
\end{yizhu}

现在我们可以陈述镜像对称猜想的几何对应物了。

\begin{conj}
对偶的 Calabi--Yau 流形族的极大退化的光滑部分，是带对偶整仿射结构的对偶
Monge--Ampère 流形。
\end{conj}

带整仿射结构的 Monge--Ampère 流形是 Calabi--Yau 流形的实类比。事实上，本节意义下的
镜像对偶对更大的一类流形也成立。我们仿照定义 2 定义一个 \emph{AK-流形}（AK 代表
\emph{affine} 与 \emph{Kähler}），但去掉条件 c)（Monge--Ampère 方程），另见 \cite{CY}。
读者容易验证，本节的所有构造（包括环面纤维化的对偶）对 AK-流形同样成立。

\begin{rmk}
把 Legendre 变换用于镜像对称的想法已存在一段时间（例如见 \cite{H}、\cite{Le}）。
\end{rmk}

\begin{rmk}
在我们对几何镜像对称的描述中忽略了 B 场。今后我们总假设 $B=0$。
\end{rmk}

\subsection{关于与非阿基米德几何之关系的猜测}

来自 CFT 与微分几何的考虑表明，$Y$ 上的整仿射结构与 Calabi--Yau 度量的凯勒类之选取无关。
于是我们得到局部域 $K=\C((q))$ 上（极大退化的）Calabi--Yau 簇的一个「组合」不变量
$(Y,T^{\Z}_Y)$。可以论证，此时存在一个典范的坐标卡图册，其转移映射属于群
\[
  SAff(n,\Z):=SL(n,\Z)\ltimes\Z^n .
\]
自然的问题是：能否不用超越方法与 Calabi--Yau 度量，纯粹用代数的方式定义并计算它？我们
期待答案是肯定的。换言之，存在一种典范的方式，把数据 $(Y,T^{\Z}_Y)$ 联系到任意光滑射影
簇 $X$（$c_1(T_X)=0$）上，只要它在任意带离散赋值的域 $K$ 上有「极大退化」。

猜想的答案（只对 $Y$ 的紧化 $\bar Y$ 而言）如下：在 $K$ 作扩域之后，选取一个具稳定约化的
模型。称特殊纤维 $X_0$ 的不可约分支 $D$ 为本质的（essential），如果 $X$ 上整体体积元在
$D$ 处的极点阶在所有 $X_0$ 的分支中达到最大。我们定义拓扑空间 $\bar Y(X_0)$ 为由本质除子
张成的 Clemens 复形（见 \cite{LTY}）。粗略地说，$\bar Y$ 的 $k$-胞腔对应于本质除子的
$(k+1)$ 重相交的不可约分支。

利用动机积分与 Berkovich 非阿基米德解析空间理论的想法（见 \cite{Be}），可以证明下述结果。

\begin{prp}
对具稳定约化的模型的不同选取，空间 $\bar Y(X_0)$ 可以典范地等同。
\end{prp}

\begin{yizhu}[Clemens 复形]
Clemens 复形是退化簇的特殊纤维 $X_0$ 所对应的一个组合模型：每个本质除子对应一个顶点，
若干个本质除子的交非空时就把相应顶点连成单形，从而得到 $\bar Y(X_0)$ 的胞腔分解。它把
$\bar Y$ 的拓扑化归为 $X_0$ 的相交理论，这正是原文把 $\bar Y(X_0)$ 称为「拓扑空间」的
含义（见 \cite{LTY}）。
\end{yizhu}

\begin{proof}
我们简述证明的想法\footnote{M.~K. 于 2000 年 10 月 30 日在 Rutgers 大学所作报告之后。}。
依照 \cite{Be}，可以对局部域 $K$ 上的代数簇 $X$ 联系一个 $K$-解析空间 $X^{an}$。对仿射的
$X$，$X^{an}$ 的点是函数代数 $\OO(X)$ 上取值于 $\R$ 的乘性半范数，它延拓 $K$ 上给定的
范数。在我们的情形 $K=\C((q))$。

\textbf{第 1 步.} 定义除子型点之集 $D(X,K)$ 为特殊纤维的不可约分支之集，对 $X$ 在
$\C[[q]]$ 上的一切模型（不必光滑）取正向极限。有自然的嵌入 $D(X,K)\subset X^{an}$：
与 $D$ 对应的半范数，由体积形式在除子 $D$ 上的极点阶取指数给出。

类似地定义 $D(X,K')=D(X\otimes K')/\mathrm{Gal}(K'/K)$，其中 $K'/K$ 是有限阶 Galois
扩域。对一切 $K'$ 取 $D(X,K')$ 的正向极限，便得到集合 $D(X)$。

\textbf{第 2 步.} 可以验证 $D(X)\subset X^{an}$，且事实上 $D(X)$ 是 $X^{an}$ 中的稠密
子集。对剩余域有限或可数的局部域，$D(X)$ 可数，而 $X^{an}$ 连通。

\textbf{第 3 步.} 给定除子 $D\in D(X)$，定义数
\[
  v(vol,D)=\frac{\operatorname{ord}_{(X\otimes K',D)}(vol)+1}
                  {\operatorname{ram}(K'/K)},
\]
其中 $\operatorname{ord}$ 表示体积形式 $vol$ 在 $D\subset X\otimes K'$ 上的零点阶，
$\operatorname{ram}$ 表示分歧指数。可以验证 $v(vol,D)$ 良定义。在 $p$ 进情形，这个数可以
由 $D$ 对积分 $\int_{X(K)}|vol|_K$ 的贡献得出，其中 $|vol|_K$ 是相伴测度。

\textbf{第 4 步.} 可以验证 $\bar Y(X_0)$ 是 $D(X)$ 中满足 $v(vol,D)$ 取极小的除子 $D$
之集合在 $X^{an}$ 中的闭包。因此它与模型的选取无关。这就解释了定理的陈述。
\end{proof}

\begin{yizhu}[「定理」应为命题 3]
原文此处作「这解释了定理的陈述」，但紧随其前陈述的是命题 3（原文 p20），本节并无相应的
定理；据上下文，此处所指即命题 3。译文照原文译出。
\end{yizhu}

在来自环面几何的例子中，空间 $\bar Y=\bar Y(X_0)$ 总是流形。目前尚不清楚光滑部分
$Y\subset\bar Y$ 以及其上的仿射结构从何而来。按猜想，这一切来自映射
$\pi:X(\bar K)\to\bar Y$，其中 $\bar K$ 是 $K$ 的代数闭包。在环面纤维化的微分几何图景中
（此时 $K=\C_{mer}$），映射 $\pi$ 是显然的：它把点的亚纯（有限分歧）族 $x_q\in X_q$
对应到度量意义下的极限点 $\lim_{q\to0}x_q\in Y$。此外，微分几何图景还提示：若
$Z\subset X$ 是代数子簇，则像 $\pi(Z(\bar K))$ 的闭包应当是 $Y$ 的分片线性闭子集，
其线性片段具有有理方向。特别地，若 $Z$ 是曲线，则 $\pi(Z(\bar K))$ 是 $Y$ 中的图像。这为
用 $Y$ 上某个量子场论的 Feynman 展开来表达 $X$ 的 Gromov--Witten 不变量开辟了道路。

\begin{yizhu}[原文的一处笔误]
原文此处写作「where $Z\subset K$ is an algebraic subvariety」，其中 $Z\subset K$ 应为
$Z\subset X$：$Z$ 是 $X$ 的代数子簇，故 $Z(\bar K)$ 落在 $\pi$ 的定义域 $X(\bar K)$ 中，
$\pi(Z(\bar K))$ 也才有意义（$K$ 是基域，不能充当子簇）。译文已按 $Z\subset X$ 译出。
\end{yizhu}

% [原文 p13–p21]
